The relay satellite revolves on a circular orbit, due to the gravitational
forces of the Earth and the Moon. Its net acceleration is: \hfill(4 p)
\begin{equation*}
\left(r + \frac{M}{m + M}R\right)\omega^2
  = \frac{GM}{(R + r)^2} + \frac{Gm}{r^2}, \tag{9.1}
\end{equation*}
where $R$ and $r$ denote the (constant) distances between Earth and Moon, and
between Moon and the satellite, respectively, while $M$ and $m$ stand for the
masses of Earth and Moon, respectively. Using Kepler's third law, the angular
velocity of the satellite, which is equal to the angular velocity of the Moon
can be written as \hfill(2 p)
\begin{equation*}
\omega^2 = \frac{4\pi^2}{P^2} = \frac{G(M + m)}{R^3}. \tag{9.2}
\end{equation*}
Substituting Eq.\ (9.2) into Eq.\ (9.1), eliminating the gravitational
constant ($G$), and introducing the dimensionless quantity $x = r/R$, we
obtain \hfill(2 p + 2 p)
\begin{gather*}
\frac{(M + m)x + M}{R^2} = \frac{M}{R^2(1 + x)^2} + \frac{m}{R^2 x^2}
  \tag{9.3} \\
(M + m)x + M = \frac{M}{(1 + x)^2} + \frac{m}{x^2}. \tag{9.4}
\end{gather*}
Assuming that $x = r/R \ll 1$ we can apply the approximation
$(1 + x)^{-2} \approx 1 - 2x$, which leads to the following expression:
\hfill(2 p)
\begin{equation*}
(M + m)x + M = M(1 - 2x) + \frac{m}{x^2}, \tag{9.5}
\end{equation*}
from which we get \hfill(2 p)
\begin{equation*}
x^3 = \frac{m}{3M + m}. \tag{9.6}
\end{equation*}
Substituting the mass ratio $m/M = 0.0123$ of the Earth--Moon system one can
obtain \hfill(2 p)
\begin{equation*}
x = 0.159\,83, \tag{9.7}
\end{equation*}
and therefore \hfill(2 p)
\begin{equation*}
r = xR_{\mathrm{Moon}} = 0.159\,83 \times 384\,400\,\mathrm{km}
  = 61\,440\,\mathrm{km} \tag{9.8}
\end{equation*}
% sic: the source writes "xR_Moon" although 384 400 km is the Earth--Moon
% distance R, not the radius of the Moon.
From this result we should subtract the radius of the Moon, and such a way we
obtain that \hfill(2 p)
\begin{equation*}
h = r - R_{\mathrm{Moon}} = 61\,440\,\mathrm{km} - 1737\,\mathrm{km}
  = \boxed{59\,703\,\mathrm{km}} \tag{9.9}
\end{equation*}

Here is another solution, which is erroneous in principle (or say,
``approximative'') but numerically gives an almost equivalent result.

One might say, that $m \ll M$ and therefore $m$ can be neglected either in
the left hand side of Eq.\ (9.1) and in the r.h.s.\ of Eq.\ (9.2), i.e.\ one
can use the following approximation:
\begin{equation*}
M + m \approx M \tag{9.10}
\end{equation*}
In this case instead of Eq.\ (9.6) one can arrive at
\begin{equation*}
x^3 = \frac{m}{3M}, \tag{9.11}
\end{equation*}
which leads to
\begin{equation*}
x = 0.160\,05, \tag{9.12}
\end{equation*}
and therefore,
\[ r = xR_{\mathrm{Moon}} = 0.160\,05 \times 384\,400\,\mathrm{km}
   = 61\,524\,\mathrm{km} \]
and, finally
\[ h = r - R_{\mathrm{Moon}} = 61\,524\,\mathrm{km} - 1737\,\mathrm{km}
   = 59\,787\,\mathrm{km} \]
This latter should be accepted as a full point solution, too. If, however,
$m$ is dropped out only from one of the two equations, then at least 2 points
should be subtracted.
